\documentclass[12pt, a4paper]{article}

% ── Paquetes ──────────────────────────────────────────────────────────────────
\usepackage[utf8]{inputenc}
\usepackage[spanish]{babel}
\usepackage{geometry}
\usepackage{graphicx}
\usepackage{listings}
\usepackage{xcolor}
\usepackage{hyperref}
\usepackage{amsmath}
\usepackage{titlesec}
\usepackage{enumitem}
\usepackage{booktabs}
\usepackage{parskip}

\geometry{margin=2.5cm}

% ── Estilo de código C++ ──────────────────────────────────────────────────────
\definecolor{codegreen}{rgb}{0.1,0.5,0.1}
\definecolor{codegray}{rgb}{0.5,0.5,0.5}
\definecolor{codeblue}{rgb}{0.1,0.1,0.7}
\definecolor{backcolour}{rgb}{0.97,0.97,0.97}

\lstdefinestyle{cpp}{
    backgroundcolor=\color{backcolour},
    commentstyle=\color{codegreen},
    keywordstyle=\color{codeblue}\bfseries,
    numberstyle=\tiny\color{codegray},
    stringstyle=\color{orange},
    basicstyle=\ttfamily\small,
    breaklines=true,
    captionpos=b,
    numbers=left,
    numbersep=5pt,
    tabsize=4,
    language=C++,
    showstringspaces=false,
    frame=single,
    rulecolor=\color{gray!30}
}
\lstset{style=cpp}

% ── Portada ───────────────────────────────────────────────────────────────────
\title{
    \vspace{2cm}
    \textbf{Entregable 2: XML} \\[0.5em]
    \large Árbol General sobre Dataset GoodReads \\[0.3em]
    \normalsize Estructura de Datos --- Semestre 1, 2026
}
\author{
    Nicolás Renato Ricciardi Acuña \\
    Enzo Gabriel Levancini Arriagada
}
\date{\today}

% ──────────────────────────────────────────────────────────────────────────────
\begin{document}

\maketitle
\tableofcontents
\newpage

% ──────────────────────────────────────────────────────────────────────────────
\section{Introducción}

Este proyecto implementa un \textbf{árbol general} en C++17 a partir de 10.000 archivos XML
del dataset GoodReads. Cada archivo describe un libro con sus metadatos (título, ISBN,
rating, año de publicación, etc.) y una lista de libros similares que define la estructura
jerárquica del árbol.

El sistema soporta tres operaciones principales sobre el árbol:
\begin{enumerate}
    \item \texttt{listar} — recorrido preorder que imprime los IDs de todos los libros.
    \item \texttt{borrar\_ratings(r)} — elimina todos los libros con rating promedio $\leq r$.
    \item \texttt{precursores(id)} — lista los libros cuyo año de publicación es menor
          al de \textit{todos} sus libros similares.
\end{enumerate}

% ──────────────────────────────────────────────────────────────────────────────
\section{Estructura del proyecto}

\begin{verbatim}
Tarea2/
├── CMakeLists.txt          <- configuración de compilación (CMake + C++17)
├── README.md
├── data/                   <- 10.000 archivos .xml del dataset GoodReads
├── libs/
│   └── pugixml/            <- librería PugiXML para parsear XML
└── src/
    ├── Book.h              <- structs Book y SimilarBook
    ├── TreeNode.h          <- nodo del árbol general
    ├── GeneralTree.h/.cpp  <- árbol general con las tres funciones requeridas
    ├── XMLParser.h/.cpp    <- carga y parseo de archivos XML
    └── main.cpp            <- punto de entrada
\end{verbatim}

% ──────────────────────────────────────────────────────────────────────────────
\section{Estructuras de datos}

\subsection{Book y SimilarBook}

\texttt{SimilarBook} almacena la información mínima de un libro referenciado en la
sección \texttt{<similar\_books>} del XML. \texttt{Book} contiene todos los campos
relevantes de un libro, incluyendo su lista de similares.

\begin{lstlisting}
struct SimilarBook {
    std::string id;
    std::string title;
    std::string isbn;
    int         publicationYear;
};

struct Book {
    std::string              id;
    std::string              title;
    std::string              isbn;
    int                      publicationYear;
    std::string              language;
    std::string              description;
    double                   averageRating;
    int                      numPages;
    std::vector<SimilarBook> similarBooks;
};
\end{lstlisting}

\subsection{TreeNode}

Cada nodo del árbol almacena un \texttt{Book} y un vector de punteros a sus hijos.
Se eligió \texttt{std::vector} sobre la representación hijo-izquierdo/hermano-derecho
para simplificar los recorridos iterativos.

\begin{lstlisting}
struct TreeNode {
    Book                   book;
    std::vector<TreeNode*> children;

    explicit TreeNode(const Book& b) : book(b) {}
};
\end{lstlisting}

\subsection{GeneralTree}

La clase principal del árbol mantiene:
\begin{itemize}
    \item \texttt{root\_} — raíz virtual con \texttt{Book} vacío, punto de entrada único.
    \item \texttt{nodeMap\_} — tabla hash \texttt{unordered\_map<string, TreeNode*>} para
          buscar cualquier nodo por ID en tiempo $O(1)$ durante la construcción.
\end{itemize}

% ──────────────────────────────────────────────────────────────────────────────
\section{Parseo de XML}

Se utiliza la librería \textbf{PugiXML} para leer los archivos XML. La estructura
real de cada archivo en el dataset es:

\begin{verbatim}
<GoodreadsResponse>
  <book>
    <id>...</id>
    <title>...</title>
    <publication_year>...</publication_year>
    <average_rating>...</average_rating>
    <similar_books>
      <book>
        <id>...</id>
        <publication_year>...</publication_year>
      </book>
    </similar_books>
  </book>
</GoodreadsResponse>
\end{verbatim}

\texttt{XMLParser::parseFile()} extrae cada campo usando \texttt{child\_value()} de
PugiXML y convierte los valores numéricos con lambdas defensivas que retornan 0
ante campos vacíos, evitando excepciones:

\begin{lstlisting}
auto toInt = [](const char* s) -> int {
    try { return (s && *s) ? std::stoi(s) : 0; }
    catch (...) { return 0; }
};
\end{lstlisting}

\texttt{XMLParser::loadDirectory()} itera recursivamente la carpeta \texttt{data/}
con \texttt{std::filesystem} y muestra una barra de progreso durante la carga.

% ──────────────────────────────────────────────────────────────────────────────
\section{Construcción del árbol}

\subsection{Problema: similitudes bidireccionales}

En el dataset GoodReads, las relaciones de similitud son \textbf{simétricas}: si el
libro A lista a B como similar, B también lista a A. Esto convierte el conjunto de
relaciones en un \textbf{grafo con ciclos}, que no puede representarse directamente
como árbol.

\subsection{Solución: DFS con conjunto de visitados}

Para convertir el grafo en un árbol sin ciclos se aplica un \textbf{DFS iterativo}
sobre todos los nodos. Cada nodo se incorpora al árbol exactamente una vez (la primera
vez que es alcanzado), garantizando que no se formen ciclos:

\begin{lstlisting}
void GeneralTree::build() {
    std::unordered_set<std::string> visited;

    for (auto& [startId, startNode] : nodeMap_) {
        if (visited.count(startId)) continue;
        visited.insert(startId);
        std::stack<TreeNode*> pila;
        pila.push(startNode);

        while (!pila.empty()) {
            TreeNode* curr = pila.top(); pila.pop();
            for (const SimilarBook& sb : curr->book.similarBooks) {
                auto it = nodeMap_.find(sb.id);
                if (it != nodeMap_.end() && !visited.count(sb.id)) {
                    visited.insert(sb.id);
                    curr->children.push_back(it->second);
                    pila.push(it->second);
                }
            }
        }
    }
    // Nodos sin padre -> hijos de la raiz virtual
    std::unordered_set<std::string> claimed;
    for (auto& [id, node] : nodeMap_)
        for (const TreeNode* hijo : node->children)
            claimed.insert(hijo->book.id);

    for (auto& [id, node] : nodeMap_)
        if (!claimed.count(id))
            root_->children.push_back(node);
}
\end{lstlisting}

Los nodos que ningún otro nodo reclama como hijo se colocan directamente bajo la
raíz virtual, formando un \textbf{bosque de expansión} con un único punto de entrada.

% ──────────────────────────────────────────────────────────────────────────────
\section{Función listar}

Recorre el árbol en \textbf{preorder} (nodo actual antes que sus hijos) e imprime
el ID de cada libro. El recorrido es \textbf{iterativo} usando una pila explícita
para evitar desbordamiento de pila con árboles profundos.

Los hijos se insertan en la pila en \textbf{orden inverso} para que al hacer
\textit{pop} se procesen de izquierda a derecha, manteniendo el orden preorder correcto.

\begin{lstlisting}
void GeneralTree::listar(std::ostream& os) const {
    std::stack<const TreeNode*> pila;
    for (auto it = root_->children.rbegin();
         it != root_->children.rend(); ++it)
        pila.push(*it);

    while (!pila.empty()) {
        const TreeNode* node = pila.top(); pila.pop();
        os << node->book.id << "\n";
        for (auto it = node->children.rbegin();
             it != node->children.rend(); ++it)
            pila.push(*it);
    }
}
\end{lstlisting}

\textbf{Complejidad:} $O(n)$, donde $n$ es el número de nodos.

% ──────────────────────────────────────────────────────────────────────────────
\section{Función borrar\_ratings}

Elimina del árbol todos los nodos cuyo \texttt{averageRating} $\leq r$. Al eliminar
un nodo se elimina también todo su subárbol (sus descendientes pierden su único
camino hacia la raíz).

\subsection{Algoritmo post-order}

Se usa un recorrido \textbf{post-order recursivo}: se evalúan los hijos \textit{antes}
que el padre. Esto permite liberar la memoria de los subárboles de abajo hacia arriba
sin perder referencias.

\begin{lstlisting}
bool GeneralTree::borrarRatingsHelper(TreeNode* node, double r) {
    std::vector<TreeNode*> sobrevivientes;
    for (TreeNode* hijo : node->children) {
        if (!borrarRatingsHelper(hijo, r))
            sobrevivientes.push_back(hijo);
        else
            deleteSubtree(hijo);
    }
    node->children = sobrevivientes;
    return node->book.averageRating <= r;
}
\end{lstlisting}

La función auxiliar retorna \texttt{true} si el nodo debe ser eliminado. El llamador
se encarga de liberarlo con \texttt{deleteSubtree()}.

\textbf{Complejidad:} $O(n)$.

% ──────────────────────────────────────────────────────────────────────────────
\section{Función precursores}

Lista los IDs de todos los libros que cumplen la siguiente condición:

\[
\forall\, s \in \text{similarBooks}(v): \quad s.\text{publicationYear} > v.\text{publicationYear}
\]

Es decir, el libro fue publicado \textit{antes} que todos sus libros similares,
y tiene al menos un libro similar.

\subsection{Manejo de años no registrados}

Si un libro similar tiene \texttt{publicationYear == 0} (campo vacío en el XML),
se trata como \textit{anterior} al libro evaluado, excluyéndolo del resultado.
Esto evita falsos positivos por datos faltantes.

\begin{lstlisting}
for (const SimilarBook& sb : sibs) {
    // año 0 = no registrado; se trata como anterior
    if (sb.publicationYear == 0 ||
        sb.publicationYear <= node->book.publicationYear) {
        todos_posteriores = false;
        break;
    }
}
\end{lstlisting}

El recorrido es \textbf{iterativo} sobre todo el árbol, usando la misma técnica de
pila explícita que \texttt{listar}.

\textbf{Complejidad:} $O(n \cdot m)$, donde $m$ es el promedio de libros similares
por nodo.

% ──────────────────────────────────────────────────────────────────────────────
\section{Estadísticas del árbol}

La función \texttt{stats()} calcula en un único recorrido $O(n)$:

\begin{table}[h]
\centering
\begin{tabular}{@{}ll@{}}
\toprule
\textbf{Métrica} & \textbf{Descripción} \\
\midrule
Nodos totales & Cantidad de libros en el árbol \\
Altura máxima & Profundidad del nodo más lejano de la raíz virtual \\
Promedio de hijos & Total de hijos / total de nodos \\
Libro más conectado & ID del nodo con mayor cantidad de hijos directos \\
\bottomrule
\end{tabular}
\caption{Métricas calculadas por \texttt{stats()}}
\end{table}

% ──────────────────────────────────────────────────────────────────────────────
\section{Compilación y ejecución}

\subsection{Requisitos}
\begin{itemize}
    \item CMake $\geq$ 3.16
    \item Compilador con soporte C++17 (GCC 8+, Clang 7+, MSVC 2017+)
\end{itemize}

\subsection{Pasos}
\begin{verbatim}
mkdir build
cd build
cmake ..
cmake --build .
./tarea2
\end{verbatim}

Los resultados se guardan automáticamente en tres archivos dentro de \texttt{build/}:
\begin{itemize}
    \item \texttt{listar.txt} — IDs en preorder
    \item \texttt{listar\_post\_borrado.txt} — IDs tras eliminar libros con rating $\leq 3.5$
    \item \texttt{precursores.txt} — IDs de libros precursores
\end{itemize}

% ──────────────────────────────────────────────────────────────────────────────
\section{Decisiones de diseño}

\begin{description}[style=nextline]
    \item[Raíz virtual]
    Se usa un nodo raíz con \texttt{Book} vacío en lugar de elegir un libro arbitrario
    como raíz real. Esto permite que el árbol sea un \textit{bosque} sin forzar una
    jerarquía artificial.

    \item[Recorridos iterativos]
    \texttt{listar} y \texttt{precursores} usan pila explícita en lugar de recursión
    para evitar desbordamiento de pila con 10.000 nodos.

    \item[DFS para romper ciclos]
    La alternativa más simple habría sido ignorar los duplicados en \texttt{children},
    pero el DFS con \texttt{visited} garantiza un árbol válido sin nodos repetidos.

    \item[Tabla hash para construcción]
    El \texttt{unordered\_map} permite buscar cualquier libro por ID en $O(1)$ durante
    \texttt{build()}, manteniendo la construcción total en $O(n \cdot m)$.
\end{description}

\end{document}
